\noindent
This document describes the three magnetic relaxation models implemented in
Fitbauer, complementary to the pre-existing singlet, doublet, sextet and
Hesse--Rübartsch distribution models:
(1)~\code{Relajacion}: a phenomenological component that interpolates between a
blocked sextet and an apparent superparamagnetic doublet;
(2)~\code{BlumeTjon}: a dynamic two-state model ($+\BHF \leftrightarrow -\BHF$)
with symmetric Bloch--McConnell-type exchange;
(3)~\code{NeelSize}: a Néel--Arrhenius model with a lognormal particle size
distribution, combining the \code{BlumeTjon} profile for each diameter.
Simultaneous global fitting of spectra at several temperatures is available from
\textbf{Fit → Global Néel (T)}.



\bigskip

%======================================================================
\section{Physical motivation}

In magnetic nanoparticles, ferrites or partially superparamagnetic iron oxides,
the magnetic moment may fluctuate during the Mössbauer observation window. If the
fluctuations are slow, the nucleus observes an almost static hyperfine field and
the spectrum is a sextet. If they are fast, the magnetic field is averaged out and
a superparamagnetic doublet or singlet appears. In the intermediate regime the
lines broaden and coalesce.

The implementation aims to avoid an automatic ``sextet + doublet = two chemical
phases'' interpretation. A doublet-like fraction may be the spectral manifestation
of fast relaxation.

\begin{aviso}
The models described here are complementary. If the \code{Relajacion} or
\code{BlumeTjon} components are not selected, the previous models are evaluated by
exactly the same code branches as before.
\end{aviso}

%======================================================================
\section{Common notation}

In the thin-absorber limit, the transmission is written as
\begin{equation}
  T(v)=b+s\,v-A(v),
\end{equation}
where $b$ is the baseline, $s$ the slope and $A(v)$ the positive absorption of the
component. For a conventional sextet:
\begin{equation}
  A_{6}(v)=D\sum_{j=1}^{6} I_j L\bigl(v; v_j,\Gamma_j\bigr),
\end{equation}
with depth $D$, intensities $I_j$ and Lorentzian profile
\begin{equation}
  L(v;v_0,\Gamma)=\frac{\Gamma^2}{(v-v_0)^2+\Gamma^2}.
\end{equation}
The first-order positions are
\begin{equation}
  v_j = \delta + q_j\,\dEQ + m_j\,\frac{\BHF}{33\,\mathrm{T}},
\end{equation}
where $m_j$ are the $\alpha$-Fe reference positions at $33$ T and $q_j$ the
first-order quadrupole pattern.

%======================================================================
\section{Phenomenological model \code{Relajacion}}

\subsection{Definition}

The \code{Relajacion} component interpolates between a blocked sextet and an
apparent superparamagnetic doublet:
\begin{equation}
  A_{\mathrm{rel}}(v)=D\left[
  f_{\mathrm{b,ef}}\,S_6(v)
  +\bigl(1-f_{\mathrm{b,ef}}\bigr)\,D_2^{\ast}(v)
  \right].
  \label{eq:rel-phen}
\end{equation}
$S_6(v)$ is the unit-depth sextet profile and $D_2^{\ast}(v)$ is the doublet
scaled so that its integrated area is comparable to that of the sextet:
\begin{equation}
  D_2^{\ast}(v)=D_2(v)\,
  \frac{\int S_6(v)\,dv}{\int D_2(v)\,dv}.
\end{equation}
In this way, $D$ approximately retains the meaning of total area/depth.

\begin{figure}[htbp]\centering
  \includegraphics[width=0.82\textwidth]{apx_relax_phenom}
  \caption{Phenomenological model \code{Relajacion} \eqref{eq:rel-phen}. As the
  blocked fraction $f_b$ decreases from $1$ (blocked sextet) to $0$ (apparent
  superparamagnetic doublet), the sextet signal is transferred to a central
  doublet of equivalent area. It does not represent two chemical phases, but the
  spectral effect of relaxation.}
  \label{fig:apx_relax_phenom}
\end{figure}

\subsection{Blocked fraction and phenomenological rate}

The user sees two specific parameters:
\begin{itemize}[leftmargin=*]
  \item $f_{\mathrm{bloq}}$: maximum blocked fraction.
  \item $\log_{10}\nu$: phenomenological relaxation rate in s$^{-1}$.
\end{itemize}
The fraction that actually enters \eqref{eq:rel-phen} is
\begin{equation}
  f_{\mathrm{b,ef}}=f_{\mathrm{bloq}}\,F(\log_{10}\nu),
  \label{eq:fb-eff}
\end{equation}
with
\begin{equation}
  F(x)=\frac{1}{1+\exp\left(\frac{x-x_c}{w}\right)},
  \qquad x_c=8.5,\quad w=0.55.
\end{equation}
Thus, slow rates ($x\ll x_c$) leave $F\simeq 1$, whereas fast rates ($x\gg x_c$)
force $F\simeq 0$.

In addition, a maximum phenomenological broadening is applied in the intermediate
regime:
\begin{equation}
  \Gamma_{\mathrm{ef}}=\Gamma\left[1+1.6\exp\left(-\frac{1}{2}
  \left(\frac{x-x_c}{0.75}\right)^2\right)\right].
\end{equation}

\begin{aviso}[Important correlation]
$f_{\mathrm{bloq}}$ and $\log_{10}\nu$ are not fully independent parameters: both
reduce the observed sextet fraction. In practice, it is advisable to fix one of
the two unless the spectrum carries enough information. To estimate fractions, fix
$\log_{10}\nu$ and free $f_{\mathrm{bloq}}$; to study dynamics, fix
$f_{\mathrm{bloq}}=1$ and free $\log_{10}\nu$.
\end{aviso}

%======================================================================
\section{Dynamic model \code{BlumeTjon}: two states}

\subsection{Exchange matrix}

The \code{BlumeTjon} component implements a first dynamic version with two opposite
field orientations:
\begin{equation}
  +\BHF \leftrightarrow -\BHF,
\end{equation}
with symmetric transition matrix
\begin{equation}
  W=\begin{pmatrix}-\nu & \nu\\ \nu & -\nu\end{pmatrix}.
\end{equation}
There is no explicit blocked fraction: the dynamics is controlled by $\nu$.

\subsection{Exchange profile for one transition}

For each sextet line, two frequencies/velocities are computed, one for $+\BHF$ and
one for $-\BHF$, denoted $v_a$ and $v_b$. A symmetric exchange form of
Bloch--McConnell/Kubo type is used. We define
\begin{equation}
  z_a = \Gamma + i(v-v_a),\qquad z_b = \Gamma + i(v-v_b),
\end{equation}
and an exchange rate in velocity units
\begin{equation}
  k_v = \frac{\nu}{\mathcal{C}_{\nu\to v}},
\end{equation}
where
\begin{equation}
  \mathcal{C}_{\nu\to v}=\frac{E_\gamma}{h c_{\mathrm{mm/s}}}.
\end{equation}
Here $E_\gamma$ is the energy of the $^{57}$Fe transition and $c_{\mathrm{mm/s}}$
the speed of light in mm/s. The absorption profile used is the real part of
\begin{equation}
  R(v)=\frac{1}{2}\,\Gamma\,
  \frac{z_a+z_b+4k_v}{z_a z_b+k_v(z_a+z_b)}.
  \label{eq:exchange}
\end{equation}
The total contribution is
\begin{equation}
  A_{\mathrm{BT}}(v)=D\sum_{j=1}^{6} I_j\,\Re\left[R_j(v)\right]_+,
\end{equation}
where $[\cdot]_+$ clips small numerical negative contributions.

\subsection{Limits}

\begin{itemize}[leftmargin=*]
  \item \textbf{Slow limit}: $k_v\to0$. Equation \eqref{eq:exchange} tends to the
  average of two static Lorentzians. In the symmetric $+\BHF/-\BHF$ pattern it
  reproduces the blocked sextet.
  \item \textbf{Intermediate}: $k_v$ comparable to the line separation. Broadening
  and coalescence appear.
  \item \textbf{Fast}: large $k_v$. The magnetic positions are averaged out and the
  spectrum tends to a doublet or singlet depending on $\dEQ$.
\end{itemize}

Figure~\ref{fig:apx_relaxation} shows the continuous evolution between these limits
as the exchange rate increases.

\begin{figure}[htbp]\centering
  \includegraphics[width=0.72\textwidth]{apx_relaxation}
  \caption{Relaxation-induced collapse in the two-state model
  $+\BHF\leftrightarrow-\BHF$ \eqref{eq:exchange}. At a low rate (bottom) the
  spectrum is the \emph{static sextet}; as the rate grows the lines broaden and
  coalesce (\emph{intermediate regime}); at a high rate (top) the magnetic
  positions are averaged to zero and a \emph{single collapsed line} remains. Curves
  offset vertically for comparison.}
  \label{fig:apx_relaxation}
\end{figure}

%======================================================================
\section{Néel--Arrhenius and size distribution}

The \code{NeelSize} component uses the Néel--Arrhenius model to convert the
particle volume into a relaxation rate. For spherical particles of diameter $d$,
\begin{equation}
  V(d)=\frac{\pi d^3}{6},
\end{equation}
and
\begin{equation}
  \tau_N(d,T)=\tau_0\exp\left(\frac{K_{\mathrm{eff}}V(d)}{k_B T}\right),
  \qquad
  \nu(d,T)=\tau_0^{-1}\exp\left(-\frac{K_{\mathrm{eff}}V(d)}{k_B T}\right).
\end{equation}
In the GUI, $\log_{10}K_{\mathrm{eff}}$ and $\log_{10}\tau_0$ are fitted to keep
reasonable numerical ranges:
\begin{equation}
  \log_{10}\nu_i=-\log_{10}\tau_0
  -\frac{K_{\mathrm{eff}}V(d_i)}{k_BT\ln 10}.
\end{equation}

The size distribution is initially modelled as lognormal. If $d_{50}$ is the
median and $\sigma_d$ the lognormal width,
\begin{equation}
  f(d)=\frac{1}{d\sigma_d\sqrt{2\pi}}
  \exp\left[-\frac{(\ln d-\ln d_{50})^2}{2\sigma_d^2}\right].
\end{equation}
The total spectrum is discretized as
\begin{equation}
  A_{\mathrm{Neel}}(v)=D\sum_i w_i\,
  A_{\mathrm{BT}}\bigl(v;\log_{10}\nu_i,T,K_{\mathrm{eff}},d_i\bigr),
\end{equation}
where $A_{\mathrm{BT}}$ is the two-state exchange profile of the previous section.
Large particles have a larger barrier $K_{\mathrm{eff}}V$ and therefore relax more
slowly; small ones relax fast and contribute a superparamagnetic signal.

The approximate blocking temperature for a given diameter is estimated with
\begin{equation}
  T_B(d)=\frac{K_{\mathrm{eff}}V(d)}{k_B\ln(\tau_M/\tau_0)},
\end{equation}
where $\tau_M$ is the Mössbauer observation time (by default $10^{-8}$ s is used
for the GUI summary).

\begin{figure}[htbp]\centering
  \includegraphics[width=0.98\textwidth]{apx_neel}
  \caption{Néel--Arrhenius: size-dependent blocking. \emph{Left}: the relaxation
  rate $\nu(d,T)$ drops as the diameter grows; small particles relax fast
  (superparamagnetic, collapsing to a singlet/doublet) and large ones remain
  \emph{blocked} (sextet). A particle appears blocked when $\nu$ falls below the
  Mössbauer observation window (grey band); as $T$ decreases the blocking shifts to
  smaller diameters. \emph{Right}: over a lognormal size distribution, at fixed $T$
  the blocking diameter $d_B$ separates the superparamagnetic fraction from the
  blocked one, and the spectrum is the mixture of both.}
  \label{fig:apx_neel}
\end{figure}

\begin{aviso}[Identifiability]
In a single spectrum, $K_{\mathrm{eff}}$, $d_{50}$, $\sigma_d$, $\tau_0$, $\Gamma$
and a possible static distribution $P(\BHF)$ may be strongly correlated. It is
therefore advisable to fix parameters using external information (TEM,
magnetometry, literature) or to use a global multi-temperature fit.
\end{aviso}

%======================================================================
\section{Global multi-temperature fit}

Fitbauer can fit several spectra at different temperatures simultaneously. For each
spectrum $r$ there is a $T_r(v)$ and a temperature $T_r$. The model shares the
physical parameters
\begin{equation}
  \theta_{\mathrm{shared}}=\{\delta,\dEQ,\BHF,\Gamma,K_{\mathrm{eff}},d_{50},
  \sigma_d,\tau_0\},
\end{equation}
and keeps local parameters per spectrum, by default
\begin{equation}
  \theta_r=\{b_r,s_r,D_r\}.
\end{equation}
The objective function is
\begin{equation}
  \chi^2=\sum_r\sum_k
  \left(\frac{T_{r,\mathrm{model}}(v_k;\theta_{\mathrm{shared}},\theta_r,T_r)
  -T_{r,\mathrm{data}}(v_k)}{\sigma_{rk}}\right)^2.
\end{equation}
This reduces the degeneracy compared to fitting a single spectrum: the same size
distribution must simultaneously explain the blocking at low $T$ and the
superparamagnetic collapse at high $T$.

%======================================================================
\section{Relation to fitting and reports}

The fittable parameters are the same as in a sextet ($\delta$, $\dEQ$, $\BHF$,
widths, depth and intensities) plus $\log_{10}\nu$ and, only in the
phenomenological model, $f_{\mathrm{bloq}}$. In \code{NeelSize}, $T$,
$\log_{10}K_{\mathrm{eff}}$, $d_{50}$, $\sigma_d$, $\log_{10}\tau_0$ and the number
of bins are added. The GUI also shows
\begin{equation}
  \tau=\frac{1}{\nu}=10^{-\log_{10}\nu}\;\mathrm{s}
\end{equation}
and, for \code{NeelSize}, an approximate $T_B$.

In reports, the \code{Relajacion} component should be interpreted as ``apparent
blocked / superparamagnetic fraction'', not as separate chemical phases. The
\code{BlumeTjon} component is reported as a dynamic two-state model and
\code{NeelSize} as a Néel--Arrhenius model with a size distribution.
